% HITSZ 实验报告模板（最终版：按用户要求的字体与字号）
% 使用 XeLaTeX 编译: xelatex report.tex
% 把校徽图片命名为 hitsz.png 并与 .tex 文件放在同一目录。

\documentclass[11pt,a4paper]{article}
\usepackage[scheme=chinese]{ctex}
\usepackage{geometry}
\geometry{left=2.4cm,right=2.4cm,top=2.2cm,bottom=2.2cm}
\usepackage{float}
\usepackage{titlesec}
\usepackage{bm}
\titleformat{\section}
  {\Large\bfseries}
  {\chinese{section}、}
  {0em}
  {}

% 字体设置 - 使用 ctex 的方式
\setmainfont{Times New Roman}
% ctex 默认已设置中文字体，如需修改可使用：
% \setCJKmainfont{SimSun}[BoldFont=SimHei]

% 预定义字号命令
\newcommand{\kaiXiaoEr}{\kaishu\fontsize{18pt}{22pt}\selectfont} % 小二
\newcommand{\heiXiaoChu}{\heiti\fontsize{36pt}{40pt}\selectfont} % 小初号
\newcommand{\heiSi}{\heiti\fontsize{14pt}{18pt}\selectfont} % 四号
\newcommand{\heiSan}{\heiti\fontsize{16pt}{20pt}\selectfont} % 三号

\usepackage{graphicx}
\usepackage{tabularx}
\usepackage{array}
\usepackage{xcolor}
\usepackage{varwidth}
\usepackage{amsmath}
\usepackage{amssymb}
\usepackage{siunitx}

\parindent=0pt
\renewcommand{\arraystretch}{1.25}

% 字段宽度
\newlength{\fieldW}
\setlength{\fieldW}{5.1cm}
\newcommand{\fieldboxW}[1][\fieldW]{\makebox[#1]{\hrulefill}}

\graphicspath{{graphic/}}

\begin{document}

% 页眉：logo 与右上文字等高显示
\begin{tabular}{@{} p{0.4\textwidth} p{0.6\textwidth} @{} }
  \raisebox{0pt}[2cm][2cm]{\includegraphics[height=2cm]{hitsz.png}} &
  \parbox[c][2cm][c]{\linewidth}{\raggedleft {\kaiXiaoEr 实验与创新实践教育中心}\\}
\\
\end{tabular}

\vspace{0.6cm}

% 标题（黑体小初号）
{\centering
  {\heiXiaoChu 实验报告}\\[6pt]
}

\vspace{0.6cm}

% 基本信息：标签与输入框均使用黑体四号
\setlength{\fieldW}{\dimexpr0.23\textwidth-2\tabcolsep\relax}
\newcommand{\fillfield}[1]{%
  \underline{%
    \makebox[\linewidth][c]{%
      \strut % 保证行高，让下划线位置统一
      % 1. 把内容存入一个临时盒子（box 0）来测量宽度
      \sbox0{#1}%
      % 2. 判断：如果内容宽度 > 列宽
      \ifdim\wd0>\linewidth
        % 则缩放到列宽
        \resizebox{\linewidth}{!}{\usebox0}%
      \else
        % 否则直接显示原内容（哪怕是空的也不会报错）
        \usebox0%
      \fi
    }%
  }%
}

\begin{tabularx}{\textwidth}{@{} 
    p{2.8cm}                  % 标签列收窄一点，留更多空间给内容
    >{\centering\arraybackslash}X 
    p{2.8cm}                  % 标签列收窄一点
    >{\centering\arraybackslash}X 
    @{}}
    
  {\heiSi 课程名称：} & {\heiSi \fillfield{电力电子系统建模与仿真}} 
    & {\heiSi 专业-班级：} & {\heiSi \fillfield{xxxx}} \\
    
  {\heiSi 实验名称：} & {\heiSi \fillfield{三相升压独立式逆变器仿真实验}} 
    & {\heiSi 学号：} & {\heiSi \fillfield{xxxxxxxxx}} \\
    
  {\heiSi 实验日期：} & {\heiSi \fillfield{2025年9月26日}} 
    & {\heiSi 姓名：} & {\heiSi \fillfield{xxx}} \\
    
  {\heiSi 同组人：} & {\heiSi \fillfield{无}} 
    & {\heiSi 报告总分数：} & {\heiSi \fillfield{}} \\ % 空内容也会有一条完整的横线
\end{tabularx}

\vspace{0.8cm}

% 蓝色分割线
{\color{blue}\rule{\textwidth}{3pt}}\\[6pt]

% 教师评语标题（黑体三号）
{\heiSan 教师评语：}\\[6pt]

% 教师评语区域：去掉外框，仅保留左侧竖线与空白区域，内容字体为黑体三号
\noindent
\begin{tabular}{@{} p{0.03\textwidth} p{0.94\textwidth} @{} }
  % 左侧细竖线
  % 留白区域（供手写评语），字体大小以黑体三号为默认（但评语通常为手写，此处用于占位）
  \vspace{0pt}\\
\end{tabular}

\vspace{8.0cm}

% 签字区域：教师签字与日期采用黑体三号，靠右对齐
\noindent\hfill
\begin{tabular}{@{} r @{} p{6cm} @{}}
  \heiSan 教师签字： & \rule{6cm}{0.4pt} \\
  \heiSan 日期： & \rule{6cm}{0.4pt} \\
\end{tabular}

\section{实验目的}
1. 熟悉三相升压独立式逆变器工作原理，探究闭环系统设计方法。

2. 利用所学的小信号建模方法，推导出三相升压独立式逆变器的动态线性模型，并学会利用 Bode
Diagram、Nyquist 等工具分析系统的特性，掌握三相升压独立式逆变器控制参数的 设计方法，并通过
仿真验证。

3. 培养学生规范的科研问题分析方法和解决思路，实现从硬件电路设计、系统建模、线性化 分析、控制
系统设计、性能指标、稳定性分析等全方位一体化的能力提升。

\section{实验设备及元器件}

\begin{table}[htbp]
\centering
\caption{实验设备及元器件}
\begin{tabularx}{\textwidth}{|c|>{\centering\arraybackslash}X|>{\centering\arraybackslash}X|>{\centering\arraybackslash}X|}
\hline
\textbf{序号} & \textbf{模块名称} & \textbf{型号} & \textbf{参数及功能说明} \\
\hline
1 & PC机 & -- & -- \\
\hline
2 & 仿真软件 & Matlab & -- \\
\hline
\end{tabularx}
\end{table}

\section{实验线路及原理}
本实验主电路为两级结构，前级为升压式转换器，后级为三相三线式全桥逆变器，前级
作用为维持直流母线电压，后级功能在控制逆变器输出电压，控制架构如图 1 所示。

\begin{figure}[H]
\centering
\includegraphics[width=0.8\textwidth]{image1.png}
\caption{三相升压独立式逆变器控制架构}
\end{figure}

\newpage
\section{主要设计参数}

\begin{table}[H]
\centering
\caption{电路设计参数}
\begin{tabular}{|c|c|}
\hline
\textbf{Item} & \textbf{Value} \\
\hline
DC Input Voltage $V_b$ & 70 V \\
\hline
DC bus Voltage $V_d$ & 100 V \\
\hline
Boost PWM $f_s$ & 40 kHz \\
\hline
Inverter PWM $f_s$ & 20 kHz \\
\hline
$C_b$ & 200 $\mu$F \\
\hline
$L_b$ & 660 $\mu$H \\
\hline
$C_{bus}$ & 940 $\mu$F \\
\hline
$L$ & 3 mH \\
\hline
$C$ & 200 $\mu$F \\
\hline
\end{tabular}
\end{table}

\section{实验内容}
\textbf{第一部分：硬件参数设计和仿真模型的调试}

1. Boost DC/DC Converter 的硬件参数设计

利用$V_d=V_b/(1-D)$，计算占空比$D$，解得$D=0.3$。令$D'=1-D=0.7$

查阅资料可知，控制-输出传递函数$G_{\nu d}\left(s\right)$为：
\begin{equation}
  G_{\nu d}\left(s\right)=\frac{V_{g}}{D^{\prime2}}\frac{1-\frac{L}{RD^{\prime2}}s}{1+\frac{L}{D^{\prime2}R}s+\frac{LC}{D^{\prime2}}s^{2}}
\end{equation}

控制-电感电流传递函数$G_{id}\left(s\right)$为：
\begin{equation}
  G_{id}(s)=\frac{V_{o}CRs+2V_{o}}{R{D^{\prime}}^{2}+CLRs^{2}+Ls}
\end{equation}
电流-电压传递函数$G_{vi}\left(s\right)$为：
\begin{equation}
  G_{vi}(s)=\frac{R{D^{\prime}}^{2}-Ls}{D^{\prime}(2+CRs)}
\end{equation}

\subsection{右半平面零点频率}
\begin{equation}
f_{rhpz} = \frac{D'^2 R}{2\pi L} = \frac{0.7^2 \times 100}{2\pi \times 6.6\times10^{-4}} \approx \SI{11.84}{kHz}
\end{equation}

\subsection{谐振频率}
\begin{equation}
f_0 = \frac{D'}{2\pi\sqrt{LC}} = \frac{0.7}{2\pi\sqrt{6.6\times10^{-4} \times 9.4\times10^{-4}}} \approx \SI{141.2}{Hz}
\end{equation}

\subsection{带宽限制原则}
\begin{itemize}
\item 电流环带宽：$f_{ci} \leq \dfrac{f_s}{10} = \SI{4}{kHz}$
\item 电压环带宽：$f_{cv} \leq \min\left(\dfrac{f_{rhpz}}{5}, \dfrac{f_0}{2}\right) \approx \min(\SI{2.37}{kHz}, \SI{70.6}{Hz})$
\end{itemize}

\subsection{电流环被控对象简化}
在电流环带宽范围内，忽略高频极点，近似为：
\begin{equation}
G_{id}(s) \approx \frac{V_d}{sL} = \frac{100}{6.6\times10^{-4}s}
\end{equation}

\subsection{设计目标}
\begin{itemize}
\item 穿越频率：$f_{ci} = \SI{3}{kHz}$
\item 相位裕度：$PM_i = 60^\circ$
\end{itemize}

\subsection{PI控制器设计}
采用标准PI控制器形式：
\begin{equation}
G_{ci}(s) = K_{pi} + \frac{K_{ii}}{s}
\end{equation}

\subsubsection{确定控制器零点位置}
设置控制器零点在穿越频率的$\dfrac{1}{3}$处：
\begin{equation}
\frac{K_{ii}}{K_{pi}} = 2\pi \times \frac{f_{ci}}{3} = 2\pi \times 1000 \approx 6283
\end{equation}

\subsubsection{计算被控对象在穿越频率处的增益}
\begin{equation}
|G_{id}(j2\pi f_{ci})| = \left|\frac{100}{j2\pi \times 3000 \times 6.6\times10^{-4}}\right| \approx \frac{100}{12.44} \approx 8.04
\end{equation}

\subsubsection{计算控制器在穿越频率处的增益}
控制器在穿越频率处的增益应满足：
\begin{equation}
|G_{ci}(j2\pi f_{ci})| \cdot |G_{id}(j2\pi f_{ci})| = 1
\end{equation}
\begin{equation}
|G_{ci}(j2\pi f_{ci})| = \frac{1}{8.04} \approx 0.124
\end{equation}

\subsubsection{计算PI参数}
PI控制器在穿越频率处的增益为：
\begin{equation}
|G_{ci}(j2\pi f_{ci})| = K_{pi} \cdot \sqrt{1 + \left(\frac{K_{ii}/K_{pi}}{2\pi f_{ci}}\right)^2}
\end{equation}

代入已知关系：
\begin{equation}
0.124 = K_{pi} \cdot \sqrt{1 + \left(\frac{6283}{2\pi \times 3000}\right)^2} = K_{pi} \cdot \sqrt{1 + (0.333)^2} = K_{pi} \cdot 1.054
\end{equation}

解得：
\begin{align}
K_{pi} &\approx \frac{0.124}{1.054} \approx 0.118 \\
K_{ii} &= K_{pi} \times 6283 \approx 0.118 \times 6283 \approx 741
\end{align}

\subsection{电流环参数调整}
基于仿真结果和实际系统响应，对参数进行微调：
\begin{itemize}
\item 最终电流环PI参数：$K_{pi} = 0.04$，$K_{ii} = 75$
\end{itemize}

\subsection{电压环被控对象}
考虑电流环闭环后的等效模型，在电压环带宽范围内，电流环可近似为单位增益：
\begin{equation}
G_{v,plant}(s) \approx G_{vd}(s)
\end{equation}

\subsection{设计目标}
\begin{itemize}
\item 穿越频率：$f_{cv} = \SI{150}{Hz}$（考虑RHP零点和谐振频率限制）
\item 相位裕度：$PM_v = 60^\circ$
\end{itemize}

\subsection{PI控制器设计}
采用标准PI控制器形式：
\begin{equation}
G_{cv}(s) = K_{pv} + \frac{K_{iv}}{s}
\end{equation}

\subsubsection{考虑RHP零点限制}
右半平面零点会引入相位滞后，限制系统带宽：
\begin{equation}
f_{cv} \leq \frac{f_{rhpz}}{3} \approx \SI{3.95}{kHz}
\end{equation}

但实际受谐振频率限制更严格：
\begin{equation}
f_{cv} \leq \frac{f_0}{2} \approx \SI{70.6}{Hz}
\end{equation}

选择$f_{cv} = \SI{150}{Hz}$作为设计目标，需要通过仿真验证稳定性。

\subsubsection{基于经验调整的参数}
由于电压环被控对象复杂，难以通过解析方法精确计算，采用基于经典控制理论和经验的参数整定方法：
\begin{itemize}
\item 初始参数：$K_{pv} = 0.05$，$K_{iv} = 50$
\item 通过仿真验证和微调
\end{itemize}

\subsection{电压环参数最终确定}
基于仿真结果和稳定性分析：
\begin{itemize}
\item 最终电压环PI参数：$K_{pv} = 0.05$，$K_{iv} = 50$
\end{itemize}

2.搭建 Boost DC/DC Converter 仿真电路模型

仿真模型如图所示：
\begin{figure}[H]
\centering
\includegraphics[width=0.8\textwidth]{image2.png}
\caption{boost电路输出电压$V_{bus}$}
\end{figure}

3.调试 Boost DC/DC Converter 的仿真电路模型，实现输出指标，检验理论设计

理论在负载为100$\Omega$时，输出电压$V_{bus}$应为100V，仿真结果如图所示：
\begin{figure}[H]
\centering
\includegraphics[width=0.8\textwidth]{image4.png}
\caption{boost电路输出电压$V_{bus}$}
\end{figure}
经测量：
$V_{bus} \approx 100V$，与理论值吻合。且输出电压纹波$\approx 0.1V$，符合设计要求。

下图为PWM波形：
\begin{figure}[H]
\centering
\includegraphics[width=0.8\textwidth]{image5.png}
\caption{PWM波形}
\end{figure}

\textbf{第二部分：系统的建模和控制设计}

1. Boost DC/DC Converter 主电路的建模及控制参数设计

见第一部分内容。

2. 三相逆变全桥逆变器主电路的搭建、坐标变换及控制参数设计

搭建三相逆变全桥逆变器仿真电路模型，见下图：
\begin{figure}[H]
\centering 
\includegraphics[width=0.8\textwidth]{image3.png}
\caption{三相逆变全桥逆变器仿真电路模型}
\end{figure}

三相光伏逆变器：通过三相全桥逆变器，受调制后 SVPWM 波控制，输出三相正弦电。
\subsection{abc 到 dq0 变换（Park 变换）}
Park 变换将三相静止坐标系 (abc) 转换为同步旋转坐标系 (dq0)，常用于电机控制和电力系统分析。

\begin{equation}
\begin{bmatrix}
x_d \\
x_q \\
x_0
\end{bmatrix}
= \frac{2}{3}
\begin{bmatrix}
\cos\theta & \cos\left(\theta - \frac{2\pi}{3}\right) & \cos\left(\theta + \frac{2\pi}{3}\right) \\
-\sin\theta & -\sin\left(\theta - \frac{2\pi}{3}\right) & -\sin\left(\theta + \frac{2\pi}{3}\right) \\
\frac{1}{2} & \frac{1}{2} & \frac{1}{2}
\end{bmatrix}
\begin{bmatrix}
x_a \\
x_b \\
x_c
\end{bmatrix}
\end{equation}
或：
\begin{align}
x_d &= \frac{2}{3} \left[ x_a \cos\theta + x_b \cos\left(\theta - \frac{2\pi}{3}\right) + x_c \cos\left(\theta + \frac{2\pi}{3}\right) \right] \\
x_q &= \frac{2}{3} \left[ -x_a \sin\theta - x_b \sin\left(\theta - \frac{2\pi}{3}\right) - x_c \sin\left(\theta + \frac{2\pi}{3}\right) \right] \\
x_0 &= \frac{1}{3} (x_a + x_b + x_c)
\end{align}
其中，$x_a, x_b, x_c$ 是三相变量，$x_d, x_q$ 是变换后的直轴和交轴变量，$x_0$ 是零序分量，$\theta$ 是旋转角度。
\subsection{dq0 到 abc 变换（逆 Park 变换）}

逆 Park 变换将同步旋转坐标系 (dq0) 转换回三相静止坐标系 (abc)。

\begin{equation}
\begin{bmatrix}
x_a \\
x_b \\
x_c
\end{bmatrix}
=
\begin{bmatrix}
\cos\theta & -\sin\theta & 1 \\
\cos\left(\theta - \frac{2\pi}{3}\right) & -\sin\left(\theta - \frac{2\pi}{3}\right) & 1 \\
\cos\left(\theta + \frac{2\pi}{3}\right) & -\sin\left(\theta + \frac{2\pi}{3}\right) & 1
\end{bmatrix}
\begin{bmatrix}
x_d \\
x_q \\
x_0
\end{bmatrix}
\end{equation}

或：
\begin{align}
x_a &= x_d \cos\theta - x_q \sin\theta + x_0 \\
x_b &= x_d \cos\left(\theta - \frac{2\pi}{3}\right) - x_q \sin\left(\theta - \frac{2\pi}{3}\right) + x_0 \\
x_c &= x_d \cos\left(\theta + \frac{2\pi}{3}\right) - x_q \sin\left(\theta + \frac{2\pi}{3}\right) + x_0
\end{align}
原理图中采集了负载侧的线电压，相电流以及角频率，通过与时间的积分得出相角偏移，作用于电压
和电流的 abc 到 dq0 的变换，送入 PI 控制器，电压外环与目标电压进行对比，输出到电流内环，最后
生成 SVPWM 波控制三相全桥逆变器。

三相逆变器在dq坐标系下：

\subsection{LC滤波器状态方程}
\begin{align}
L_f \frac{di_d}{dt} &= v_{id} - v_{od} + \omega L_f i_q \\
L_f \frac{di_q}{dt} &= v_{iq} - v_{oq} - \omega L_f i_d \\
C_f \frac{dv_{od}}{dt} &= i_d - i_{od} + \omega C_f v_{oq} \\
C_f \frac{dv_{oq}}{dt} &= i_q - i_{oq} - \omega C_f v_{od}
\end{align}

\subsection{输出方程}
\begin{align}
v_{od} &= \text{输出电压d轴分量} \\
v_{oq} &= \text{输出电压q轴分量} \\
i_{od} &= \frac{v_{od}}{R_{load}} \quad (\text{电阻负载}) \\
i_{oq} &= \frac{v_{oq}}{R_{load}}
\end{align}

\subsection{控制目标方程}
输出电压幅值：
\begin{equation}
V_{out} = \sqrt{v_{od}^2 + v_{oq}^2}
\end{equation}

期望设定：$v_{oq}^* = 0$（q轴电压为零，实现单位功率因数）

电流内环控制器设计过程如下：

\subsection{电流环被控对象}
忽略交叉耦合项，dq轴电流环可视为两个独立的一阶系统：
\begin{align}
G_{id}(s) &= \frac{1}{L_f s + R_L} \approx \frac{1}{L_f s} \\
G_{iq}(s) &= \frac{1}{L_f s + R_L} \approx \frac{1}{L_f s}
\end{align}
其中$R_L$为电感寄生电阻。

\subsection{考虑延时环节}
PWM和采样造成的总延时：
\begin{equation}
T_d = T_s + T_{PWM} = \frac{1}{2f_{sw}} + \frac{1}{2f_{sw}} = \frac{1}{f_{sw}} = \SI{50}{\micro s}
\end{equation}

延时环节传递函数：
\begin{equation}
G_d(s) = e^{-sT_d} \approx \frac{1}{1 + sT_d + \frac{(sT_d)^2}{2}}
\end{equation}

\subsection{PI控制器设计}
电流环PI控制器：
\begin{equation}
G_{ci}(s) = K_{pi} + \frac{K_{ii}}{s}
\end{equation}

\subsubsection{基于内模控制的参数整定}
\begin{align}
K_{pi} &= \frac{L_f}{2T_d} = \frac{3\times10^{-3}}{2\times50\times10^{-6}} = 30 \\
K_{ii} &= \frac{R_L}{2T_d} \approx \frac{0.1}{2\times50\times10^{-6}} = 1000 \quad (\text{假设 } R_L = {0.1}{\Omega})
\end{align}

\subsubsection{基于带宽设计的参数整定}
设电流环带宽 $f_{ci} = \SI{2}{kHz}$：
\begin{align}
K_{pi} &= 2\pi f_{ci} L_f = 2\pi\times2000\times3\times10^{-3} \approx 37.7 \\
K_{ii} &= 2\pi f_{ci} R_L \approx 2\pi\times2000\times0.1 \approx 1257
\end{align}

\subsection{前馈解耦补偿}
\begin{align}
v_{id}^* &= v_{id}^{PI} - \omega L_f i_q + v_{od} \\
v_{iq}^* &= v_{iq}^{PI} + \omega L_f i_d + v_{oq}
\end{align}

\subsection{推荐的电流环参数}
\begin{align}
K_{pi} &= 35 \\
K_{ii} &= 1200
\end{align}

电压外环控制器设计过程如下：

\subsection{电压环被控对象}
电流环闭环后可近似为单位增益，电压环被控对象：
\begin{equation}
G_v(s) = \frac{1}{C_f s} \cdot \frac{1}{1 + sT_d}
\end{equation}

\subsection{PR控制器设计}
为消除稳态误差，采用PR控制器：
\begin{equation}
G_{cv}(s) = K_{pv} + \frac{2K_{rv}\omega_c s}{s^2 + 2\omega_c s + \omega_0^2}
\end{equation}
其中$\omega_0 = 2\pi\times50 = \SI{314}{rad/s}$，$\omega_c$为谐振带宽。

\subsection{基于带宽设计的参数整定}
设电压环带宽 $f_{cv} = \SI{200}{Hz}$：
\begin{align}
K_{pv} &= 2\pi f_{cv} C_f = 2\pi\times200\times200\times10^{-6} \approx 0.251 \\
K_{rv} &= 10 \sim 20 \quad (\text{根据谐波抑制要求调整}) \\
\omega_c &= 2\pi\times5 = \SI{31.4}{rad/s} \quad (\text{谐振带宽})
\end{align}

\subsection{前馈补偿}
\begin{align}
i_d^* &= i_d^{PR} + \omega C_f v_{oq} + i_{od} \\
i_q^* &= i_q^{PR} - \omega C_f v_{od} + i_{oq}
\end{align}

\subsection{推荐的电压环参数}
\begin{align}
K_{pv} &= 0.25 \\
K_{rv} &= 15 \\
\omega_c &= \SI{31.4}{rad/s}
\end{align}

SVPWM实现与调制策略如下：

\subsection{空间矢量定义}
参考电压矢量：
\begin{equation}
V_{ref} = v_\alpha + j v_\beta = \frac{2}{3}(v_a + \alpha v_b + \alpha^2 v_c), \quad \alpha = e^{j\frac{2\pi}{3}}
\end{equation}

\subsection{扇区判断}
\begin{equation}
\begin{cases}
v_{ref1} = v_\beta \\
v_{ref2} = \frac{\sqrt{3}}{2}v_\alpha - \frac{1}{2}v_\beta \\
v_{ref3} = -\frac{\sqrt{3}}{2}v_\alpha - \frac{1}{2}v_\beta
\end{cases}
\end{equation}

扇区号 $N$ 由$v_{ref1}, v_{ref2}, v_{ref3}$的符号决定。

\subsection{作用时间计算}
\begin{align}
X &= \frac{\sqrt{3}v_\beta T_s}{V_{dc}} \\
Y &= \frac{\frac{3}{2}v_\alpha + \frac{\sqrt{3}}{2}v_\beta}{V_{dc}}T_s \\
Z &= \frac{-\frac{3}{2}v_\alpha + \frac{\sqrt{3}}{2}v_\beta}{V_{dc}}T_s
\end{align}

\subsection{调制深度限制}
最大线性调制范围：
\begin{equation}
m_{max} = \frac{V_{dc}}{\sqrt{3}V_{out,max}} = \frac{100}{\sqrt{3}\times70} \approx 0.825
\end{equation}

过调制策略：当$m > 0.825$时采用过调制算法。

3. 测试不同工况下的系统性能。

工况一：$R_{load} = 10\Omega$，$V_{o}=50V$
仿真结果如图所示：
\begin{figure}[H]
\centering
\includegraphics[width=0.8\textwidth]{image6.png}
\caption{$R_{load} = 10\Omega$，$V_{o}=50V$输出电压波形}
\end{figure}
输出稳定时测量：
$V_{od} \approx 49.94V$，$V_{oq} \approx 0V$，与理论值吻合。符合设计要求。
\begin{figure}[H]
\centering
\includegraphics[width=0.8\textwidth]{image17.png}
\caption{$R_{load} = 10\Omega$，$V_{o}=50V$输出电流波形}
\end{figure}
输出稳定时测量：
$I_{od} \approx 6.428A$，$I_{oq} \approx 0A$。

工况二：$R_{load} = 50\Omega$，$V_{o}=50V$
仿真结果如图所示：
\begin{figure}[H]
\centering
\includegraphics[width=0.8\textwidth]{image7.png}
\caption{$R_{load} = 50\Omega$，$V_{o}=50V$输出电压波形}
\end{figure}
输出稳定时测量：
$V_{od} \approx 50.48V$，$V_{oq} \approx 0V$，与理论值吻合。符合设计要求。
\begin{figure}[H]
\centering
\includegraphics[width=0.8\textwidth]{image16.png}
\caption{$R_{load} = 50\Omega$，$V_{o}=50V$输出电流波形}
\end{figure}
输出稳定时测量：
$I_{od} \approx 4.086A$，$I_{oq} \approx 0A$。

工况三：$R_{load} = 10\Omega$，$V_{o}=60V$
仿真结果如图所示：
\begin{figure}[H]
\centering
\includegraphics[width=0.8\textwidth]{image8.png}
\caption{$R_{load} = 10\Omega$，$V_{o}=60V$输出电压波形}
\end{figure}
输出稳定时测量：
$V_{od} \approx 59.70V$，$V_{oq} \approx 0V$，与理论值吻合。符合设计要求。
\begin{figure}[H]
\centering
\includegraphics[width=0.8\textwidth]{image14.png}
\caption{$R_{load} = 10\Omega$，$V_{o}=60V$输出电流波形}
\end{figure}
输出稳定时测量：
$I_{od} \approx 7.531A$，$I_{oq} \approx 0A$。

工况四：$R_{load} = 50\Omega$，$V_{o}=60V$
仿真结果如图所示：
\begin{figure}[H]
\centering
\includegraphics[width=0.8\textwidth]{image9.png}
\caption{$R_{load} = 50\Omega$，$V_{o}=60V$输出电压波形}
\end{figure}
输出稳定时测量：
$V_{od} \approx 60.50V$，$V_{oq} \approx 0V$，与理论值吻合。符合设计要求。
\begin{figure}[H]
\centering
\includegraphics[width=0.8\textwidth]{image15.png}
\caption{$R_{load} = 50\Omega$，$V_{o}=60V$输出电流波形}
\end{figure}
输出稳定时测量：
$I_{od} \approx 4.725A$，$I_{oq} \approx 0A$。

工况五：$R_{load} = 10\Omega$，$V_{o}=70V$
仿真结果如图所示：
\begin{figure}[H]
\centering
\includegraphics[width=0.8\textwidth]{image10.png}
\caption{$R_{load} = 10\Omega$，$V_{o}=70V$输出电压波形}
\end{figure}
输出稳定时测量：
$V_{od} \approx 62.24V$，$V_{oq} \approx 0V$，这是由于过调制，输出电压已低于设定值。
\begin{figure}[H]
\centering
\includegraphics[width=0.8\textwidth]{image13.png}
\caption{$R_{load} = 10\Omega$，$V_{o}=70V$输出电流波形}
\end{figure}
输出稳定时测量：
$I_{od} \approx 7.762A$，$I_{oq} \approx 0A$。

工况六：$R_{load} = 50\Omega$，$V_{o}=70V$
仿真结果如图所示：
\begin{figure}[H]
\centering
\includegraphics[width=0.8\textwidth]{image11.png}
\caption{$R_{load} = 50\Omega$，$V_{o}=70V$输出电压波形}
\end{figure}
输出稳定时测量：
$V_{od} \approx 64.53V$，$V_{oq} \approx 0V$，这是由于过调制，输出电压已低于设定值。
\begin{figure}[H]
\centering
\includegraphics[width=0.8\textwidth]{image12.png}
\caption{$R_{load} = 50\Omega$，$V_{o}=70V$输出电流波形}
\end{figure}
输出稳定时测量：
$I_{od} \approx 4.817A$，$I_{oq} \approx 0A$。

输出电压一定时，当负载增大，输出电流均会有所上升，符合预期。
\end{document}
